% Options for packages loaded elsewhere
% Options for packages loaded elsewhere
\PassOptionsToPackage{unicode}{hyperref}
\PassOptionsToPackage{hyphens}{url}
\PassOptionsToPackage{dvipsnames,svgnames,x11names}{xcolor}
%
\documentclass[
  letterpaper,
  DIV=11,
  numbers=noendperiod]{scrartcl}
\usepackage{xcolor}
\usepackage{amsmath,amssymb}
\setcounter{secnumdepth}{-\maxdimen} % remove section numbering
\usepackage{iftex}
\ifPDFTeX
  \usepackage[T1]{fontenc}
  \usepackage[utf8]{inputenc}
  \usepackage{textcomp} % provide euro and other symbols
\else % if luatex or xetex
  \usepackage{unicode-math} % this also loads fontspec
  \defaultfontfeatures{Scale=MatchLowercase}
  \defaultfontfeatures[\rmfamily]{Ligatures=TeX,Scale=1}
\fi
\usepackage{lmodern}
\ifPDFTeX\else
  % xetex/luatex font selection
\fi
% Use upquote if available, for straight quotes in verbatim environments
\IfFileExists{upquote.sty}{\usepackage{upquote}}{}
\IfFileExists{microtype.sty}{% use microtype if available
  \usepackage[]{microtype}
  \UseMicrotypeSet[protrusion]{basicmath} % disable protrusion for tt fonts
}{}
\makeatletter
\@ifundefined{KOMAClassName}{% if non-KOMA class
  \IfFileExists{parskip.sty}{%
    \usepackage{parskip}
  }{% else
    \setlength{\parindent}{0pt}
    \setlength{\parskip}{6pt plus 2pt minus 1pt}}
}{% if KOMA class
  \KOMAoptions{parskip=half}}
\makeatother
% Make \paragraph and \subparagraph free-standing
\makeatletter
\ifx\paragraph\undefined\else
  \let\oldparagraph\paragraph
  \renewcommand{\paragraph}{
    \@ifstar
      \xxxParagraphStar
      \xxxParagraphNoStar
  }
  \newcommand{\xxxParagraphStar}[1]{\oldparagraph*{#1}\mbox{}}
  \newcommand{\xxxParagraphNoStar}[1]{\oldparagraph{#1}\mbox{}}
\fi
\ifx\subparagraph\undefined\else
  \let\oldsubparagraph\subparagraph
  \renewcommand{\subparagraph}{
    \@ifstar
      \xxxSubParagraphStar
      \xxxSubParagraphNoStar
  }
  \newcommand{\xxxSubParagraphStar}[1]{\oldsubparagraph*{#1}\mbox{}}
  \newcommand{\xxxSubParagraphNoStar}[1]{\oldsubparagraph{#1}\mbox{}}
\fi
\makeatother


\usepackage{longtable,booktabs,array}
\usepackage{calc} % for calculating minipage widths
% Correct order of tables after \paragraph or \subparagraph
\usepackage{etoolbox}
\makeatletter
\patchcmd\longtable{\par}{\if@noskipsec\mbox{}\fi\par}{}{}
\makeatother
% Allow footnotes in longtable head/foot
\IfFileExists{footnotehyper.sty}{\usepackage{footnotehyper}}{\usepackage{footnote}}
\makesavenoteenv{longtable}
\usepackage{graphicx}
\makeatletter
\newsavebox\pandoc@box
\newcommand*\pandocbounded[1]{% scales image to fit in text height/width
  \sbox\pandoc@box{#1}%
  \Gscale@div\@tempa{\textheight}{\dimexpr\ht\pandoc@box+\dp\pandoc@box\relax}%
  \Gscale@div\@tempb{\linewidth}{\wd\pandoc@box}%
  \ifdim\@tempb\p@<\@tempa\p@\let\@tempa\@tempb\fi% select the smaller of both
  \ifdim\@tempa\p@<\p@\scalebox{\@tempa}{\usebox\pandoc@box}%
  \else\usebox{\pandoc@box}%
  \fi%
}
% Set default figure placement to htbp
\def\fps@figure{htbp}
\makeatother





\setlength{\emergencystretch}{3em} % prevent overfull lines

\providecommand{\tightlist}{%
  \setlength{\itemsep}{0pt}\setlength{\parskip}{0pt}}



 


\KOMAoption{captions}{tableheading}
\makeatletter
\@ifpackageloaded{caption}{}{\usepackage{caption}}
\AtBeginDocument{%
\ifdefined\contentsname
  \renewcommand*\contentsname{Table of contents}
\else
  \newcommand\contentsname{Table of contents}
\fi
\ifdefined\listfigurename
  \renewcommand*\listfigurename{List of Figures}
\else
  \newcommand\listfigurename{List of Figures}
\fi
\ifdefined\listtablename
  \renewcommand*\listtablename{List of Tables}
\else
  \newcommand\listtablename{List of Tables}
\fi
\ifdefined\figurename
  \renewcommand*\figurename{Figure}
\else
  \newcommand\figurename{Figure}
\fi
\ifdefined\tablename
  \renewcommand*\tablename{Table}
\else
  \newcommand\tablename{Table}
\fi
}
\@ifpackageloaded{float}{}{\usepackage{float}}
\floatstyle{ruled}
\@ifundefined{c@chapter}{\newfloat{codelisting}{h}{lop}}{\newfloat{codelisting}{h}{lop}[chapter]}
\floatname{codelisting}{Listing}
\newcommand*\listoflistings{\listof{codelisting}{List of Listings}}
\makeatother
\makeatletter
\makeatother
\makeatletter
\@ifpackageloaded{caption}{}{\usepackage{caption}}
\@ifpackageloaded{subcaption}{}{\usepackage{subcaption}}
\makeatother
\usepackage{bookmark}
\IfFileExists{xurl.sty}{\usepackage{xurl}}{} % add URL line breaks if available
\urlstyle{same}
\hypersetup{
  pdftitle={Problem-Solving and Critical Path Reasoning (PSCPR)},
  colorlinks=true,
  linkcolor={blue},
  filecolor={Maroon},
  citecolor={Blue},
  urlcolor={Blue},
  pdfcreator={LaTeX via pandoc}}


\title{Problem-Solving and Critical Path Reasoning (PSCPR)}
\author{}
\date{}
\begin{document}
\maketitle


\section{Problem-Solving and Critical Path Reasoning
(PSCPR)}\label{problem-solving-and-critical-path-reasoning-pscpr}

Download PDF Download PDF (with examples)

Effective problem-solving requires consciously transitioning from
comfortable certainty into productive uncertainty. Initially, we rely on
clear facts and well-established methods. But genuine insight demands
stepping beyond what's known to recognize gaps, implicit assumptions,
and hidden perspectives. True mastery emerges when we courageously
embrace exploration, engaging the unknown and transforming uncertainty
into discovery and deeper understanding.

\textbf{Certainty → Uncertainty → Exploration}

\subsection{1) Observation (Known Knowns --- Type 1
Expert)}\label{observation-known-knowns-type-1-expert}

\begin{itemize}
\tightlist
\item
  \textbf{Questions:} ``What is given? What do I know? What did I know
  before?''
\item
  \textbf{Givens:} Measurements, observations, quantities, data, or
  explicitly stated conditions
\item
  \textbf{Knowns:} Concepts, theories, principles, relationships, or
  frameworks
\item
  \textbf{Visuals:} Graphs, charts, diagrams. Sketch out thoughts and
  ideas
\end{itemize}

\subsection{2) Analysis (Known Unknowns --- Type 1
Expert)}\label{analysis-known-unknowns-type-1-expert}

\begin{itemize}
\tightlist
\item
  \textbf{Questions:} ``What am I trying to find? What do I know I don't
  know? What do I need to know?''
\item
  \textbf{Objectives:} Identify desired values, parameters, relations,
  or outcomes
\item
  \textbf{Unknowns:} Consider quantities, parameters, or relationships
  required to obtain objectives
\end{itemize}

\subsection{3) Inference (Unknown Knowns --- Type 2
Expert)}\label{inference-unknown-knowns-type-2-expert}

\begin{itemize}
\tightlist
\item
  \textbf{Questions:} ``What am I missing? What do I not know that I
  know? Has the problem been sufficiently reduced? Do I know/have enough
  to determine what needs to be known? What am I not acknowledging? What
  can be deduced?''
\item
  \textbf{Alternatives:} Consider multiple perspectives, challenge
  hidden assumptions, acknowledge implicit or overlooked knowledge
  (revisit Step 1).
\item
  \textbf{Complexities:} Reduce degrees of freedom, simplify the problem
  wherever possible (but no further).
\item
  \textbf{Subtleties:} Recognize possible equations, relations, or
  quantities that aren't explicitly given yet are implicitly known and
  could be leveraged to solve the problem.
\end{itemize}

\subsection{4) Exploration (Unknown Unknowns --- Type 3
Expert)}\label{exploration-unknown-unknowns-type-3-expert}

\begin{itemize}
\tightlist
\item
  \textbf{Questions:} ``Does this make sense? What do I not know that I
  do not know? What am I unaware of or not perceiving? What haven't I
  considered? What must be induced?''
\item
  \textbf{Curiosities:} Explore original ideas, embrace uncertainties.
  Perform tests and experiments, run simulations, gather observations.

  \begin{itemize}
  \tightlist
  \item
    You don't (won't) know until you look.
  \item
    ``The cave you fear to enter holds the treasure you seek.'' ---
    Joseph Campbell
  \end{itemize}
\item
  \textbf{Answers:} Perform sanity checks, redundancy checks, and
  consistency checks.

  \begin{itemize}
  \tightlist
  \item
    Dimensional analysis
  \end{itemize}
\item
  \textbf{Reflections:} Recalculate, reformulate, or revise approach as
  necessary.
\item
  \textbf{Refinements:} Iteratively adjust methodologies, hypotheses,
  and solutions.
\end{itemize}

\begin{center}\rule{0.5\linewidth}{0.5pt}\end{center}

\subsection{How to Use the PSCPR Guide
Sheet}\label{how-to-use-the-pscpr-guide-sheet}

\subsubsection{1. Observation --- Establish Claim, Null, and
Assumptions}\label{observation-establish-claim-null-and-assumptions}

\begin{itemize}
\tightlist
\item
  In the \textbf{Claim} line, write the statement under test (the
  proposed explanation or hypothesis).
\item
  In the \textbf{Null} line, write the ``business-as-usual'' alternative
  that holds if (P) is not true.
\item
  In \textbf{Assumptions}, list the key conditions, constants, knowns,
  givens, or context considered.
\end{itemize}

\subsubsection{2. Analysis --- Define Necessary
Observables}\label{analysis-define-necessary-observables}

\begin{itemize}
\tightlist
\item
  Under ``If (P) is true\ldots{}'', record one or more \textbf{Necessary
  Observables (Questions)} that must hold if the Claim stands.
\item
  Under ``If (N) is true\ldots{}'', record one or more \textbf{Necessary
  Observables (Questions)} that must hold if the Null stands.
\item
  A \textbf{Necessary Observable} is a condition whose clear failure
  would directly undermine that position, turning each position into
  testable expectations.
\end{itemize}

\subsubsection{3. Inference --- Organize evidence and candidate
stories}\label{inference-organize-evidence-and-candidate-stories}

\begin{itemize}
\tightlist
\item
  List concrete \textbf{facts} known to be true in this specific case
  (measurements, observations, calculations, documented events).
\item
  For each fact, mark whether it falsifies the Claim, the Null, both, or
  neither, based on the Necessary Observables.
\item
  List relevant \textbf{patterns} (how similar systems usually behave,
  base-rates, historical tendencies).
\item
  For each pattern, mark whether it rejects or weakens the Claim, the
  Null, both, or neither.
\item
  In the first story line, describe the behavior if the Claim were true
  (\textbf{Claim story}).
\item
  In the second story line, describe the behavior if the Null were true
  (\textbf{Null story}).
\item
  Indicate which story currently fits the recorded Facts and Patterns
  with the fewest added assumptions (Claim story, Null story, tie, or
  neither).
\item
  Stories in this section function as candidate models; Facts and
  Patterns determine which models remain viable.
\end{itemize}

\subsubsection{4. Exploration --- Apply falsification and test
outcome}\label{exploration-apply-falsification-and-test-outcome}

\begin{itemize}
\tightlist
\item
  Answer the falsification prompts to indicate whether any Fact or
  Pattern clearly contradicts a Necessary Observable for the Claim or
  for the Null.
\item
  Identify the Outcome and select an appropriate Claim Status.
\end{itemize}

\subsubsection{5. Hypothesis}\label{hypothesis}

Write a brief statement summarizing the current position. Before writing
the hypothesis, pause and ask:

\begin{itemize}
\tightlist
\item
  Does it follow from the facts ((D)), patterns ((I)), and the questions
  (Q\_P / Q\_N)? Was anything ignored that clearly hits (P) or (N)?
\item
  Is there a reasonable rival story that still explains (D + I) almost
  as well (and should that become a new (P) or (N) next round)?
\item
  Is anything in (A) really an untested claim that needs its own (P/N)
  check?
\item
  And what single new observation or test would most change my rating
  for (P)?
\end{itemize}

\begin{center}\rule{0.5\linewidth}{0.5pt}\end{center}

\subsection{PSCPR Guide Sheet
(fillable)}\label{pscpr-guide-sheet-fillable}

\begin{quote}
\textbf{Test a claim against a null using facts (deduction), patterns
(induction), and stories (abduction).}
\end{quote}

\subsubsection{Observation}\label{observation}

\begin{itemize}
\tightlist
\item
  \textbf{Claim (P)} (paradigm):\\
  \emph{\ldots{}}
\item
  \textbf{Null (N)} (status quo, business-as-usual):\\
  \emph{\ldots{}}
\item
  \textbf{Assumptions (A)} (context/givens):\\
  \emph{\ldots{}}
\end{itemize}

\subsubsection{Analysis}\label{analysis}

If (P) is true, then the following are \textbf{Necessary Observables} of
(P) (Questions):

\begin{itemize}
\tightlist
\item
  (Q\_\{P1\}): \emph{\ldots{}}
\item
  (Q\_\{P2\}): \emph{\ldots{}}
\end{itemize}

If (N) is true, then the following are \textbf{Necessary Observables} of
(N) (Questions):

\begin{itemize}
\tightlist
\item
  (Q\_\{N1\}): \emph{\ldots{}}
\item
  (Q\_\{N2\}): \emph{\ldots{}}
\end{itemize}

\subsubsection{Inference}\label{inference}

\textbf{Facts (Deduction) --- What is true?} (observation / measurement)

\begin{itemize}
\tightlist
\item
  (D\_1): \emph{\ldots{}}\\
  This fact falsifies:

  \begin{itemize}
  \tightlist
  \item[$\square$]
    (Q\_P)\\
  \item[$\square$]
    (Q\_N)\\
  \item[$\square$]
    Both\\
  \item[$\square$]
    Neither
  \end{itemize}
\item
  (D\_2): \emph{\ldots{}}\\
  This fact falsifies:

  \begin{itemize}
  \tightlist
  \item[$\square$]
    (Q\_P)\\
  \item[$\square$]
    (Q\_N)\\
  \item[$\square$]
    Both\\
  \item[$\square$]
    Neither
  \end{itemize}
\end{itemize}

\textbf{Patterns (Induction) --- What usually happens?} (base-rate /
prior behavior)

\begin{itemize}
\tightlist
\item
  (I\_1): \emph{\ldots{}}\\
  This pattern rejects:

  \begin{itemize}
  \tightlist
  \item[$\square$]
    (Q\_P)\\
  \item[$\square$]
    (Q\_N)\\
  \item[$\square$]
    Both\\
  \item[$\square$]
    Neither
  \end{itemize}
\item
  (I\_2): \emph{\ldots{}}\\
  This pattern rejects:

  \begin{itemize}
  \tightlist
  \item[$\square$]
    (Q\_P)\\
  \item[$\square$]
    (Q\_N)\\
  \item[$\square$]
    Both\\
  \item[$\square$]
    Neither
  \end{itemize}
\end{itemize}

\textbf{Stories (Abduction) --- What possibly happens?} (If (P)/(N) is
true, how/why?)

\begin{itemize}
\tightlist
\item
  (S\_P): \emph{\ldots{}}\\
\item
  (S\_N): \emph{\ldots{}}
\end{itemize}

Given the Facts and Patterns, which Story fits better with fewer
assumptions?

\begin{itemize}
\tightlist
\item[$\square$]
  (S\_P)
\item[$\square$]
  (S\_N)
\item[$\square$]
  Tie
\item[$\square$]
  Neither
\end{itemize}

\begin{quote}
\textbf{Note:} Stories are to be tested, not believed.
\end{quote}

\subsubsection{Exploration}\label{exploration}

Did any fact or pattern falsify (P) (a necessary (Q\_P) clearly fails)?
If so, which?

\begin{itemize}
\tightlist
\item[$\square$]
  Yes
\item[$\square$]
  No\\
  Which: \emph{\ldots{}}
\end{itemize}

Did any fact or pattern falsify (N) (a necessary (Q\_N) clearly fails)?
If so, which?

\begin{itemize}
\tightlist
\item[$\square$]
  Yes
\item[$\square$]
  No\\
  Which: \emph{\ldots{}}
\end{itemize}

\textbf{Outcome}

\begin{itemize}
\tightlist
\item[$\square$]
  \textbf{Retain (N)} --- Evidence is too weak to dismiss the null; (P)
  is not yet better than (N).
\item[$\square$]
  \textbf{Accept (P)} --- Provisionally accept (P); (N) is not
  falsified, but (P) is better than (N).
\item[$\square$]
  \textbf{Accept (P), Reject (N)} --- (N) is falsified by a failed
  (Q\_N) or clear conflict with (D); (P) survives.
\item[$\square$]
  \textbf{Undecided} --- (D) and (I) are insufficient or conflicting;
  need new observations or tests.
\end{itemize}

\textbf{Claim Status}

\begin{itemize}
\tightlist
\item[$\square$]
  \textbf{0.0 --- False} --- Contradicted by facts or necessary
  conditions; (P) fails (Q\_P).
\item[$\square$]
  \textbf{0.2 --- Speculative} --- Mostly a story; little or no support
  from (D) or (I), but not yet ruled out.
\item[$\square$]
  \textbf{0.4 --- Plausible} --- Consistent with (D) and (I), worth
  considering; rivals are at least as strong.
\item[$\square$]
  \textbf{0.6 --- Probable} --- Fits (D) and (I) better than (N) and
  other rivals; no serious contradictions.
\item[$\square$]
  \textbf{0.8 --- Corroborative} --- Strong fit to (D) and (I); (P)
  survived several tests; rivals are weaker.
\item[$\square$]
  \textbf{1.0 --- True} --- Operationally treated as true; no rival
  explains (D) and (I) better.
\end{itemize}

\subsubsection{Hypothesis}\label{hypothesis-1}

State your working hypothesis below, and check: does it follow from your
facts ((D)), patterns ((I)), and necessary questions ((Q\_P/Q\_N))? Is
there a live rival story left? What one test would most change your
rating for (P)?

\emph{\ldots{}}

Example: ``Given ({[}A{]}) and based on the current facts ({[}D{]}) and
patterns ({[}I{]}), ({[}P/N{]}) is the best description, so {[}retain
(N) / provisionally accept (P) / accept (P) and reject (N) / stay
undecided{]}.''

\begin{center}\rule{0.5\linewidth}{0.5pt}\end{center}

\subsection{Printable layout (original
sheet)}\label{printable-layout-original-sheet}

\begin{quote}
If you want the exact printable formatting, embed the original sheet
image (or link the PDF download above).
\end{quote}

\begin{figure}[H]

{\centering \pandocbounded{\includegraphics[keepaspectratio]{pscpr-guide-sheet.png}}

}

\caption{PSCPR Guide Sheet (print layout)}

\end{figure}%

\begin{center}\rule{0.5\linewidth}{0.5pt}\end{center}

\subsection{References}\label{references}

\begin{itemize}
\tightlist
\item
  The Ethical Skeptic, ``The Elements of Hypothesis''; The Ethical
  Skeptic, WordPress, 4 Mar 2019; Web: https://wp.me/p17q0e-94J\\
\item
  The Ethical Skeptic, ``The Three Types of Expert''; The Ethical
  Skeptic, WordPress, 28 Nov 2021; Web:
  https://theethicalskeptic.com/?p=57222
\end{itemize}




\end{document}
